% Lösungen Einheit 1 von 8 – Ereignisse als Mengen (zusammenbau v0.1)
\einheitenkopf{Einheit 1 von 8 · Ereignisse als Mengen}

\erg{12}{a) einzelnes Ergebnis \quad b) 8 Ergebnisse: (K; K; K), (K; K; Z), (K; Z; K), (K; Z; Z), (Z; K; K), (Z; K; Z), (Z; Z; K), (Z; Z; Z); Laplace: ja, alle gleich wahrscheinlich \quad c) $A \cap B$: ZKK; weder A noch B: KZZ \quad d) 3 Ergebnisse: (2; 3), (3; 2), (4; 1) \quad e) $P(S \cap H)$: Anteil aller Schüler, die Rad fahren und einen Helm tragen; $P_S(H)$: Anteil der Helmträger unter den Radfahrern \quad f) weder Hund noch jünger als zwei Jahre: $P = 1 - (0{,}45 + 0{,}3 - 0{,}1) = 0{,}35$ \quad g) $0{,}45 + 0{,}35 - 2 \cdot 0{,}15 = 0{,}5$ \quad h) Wahrscheinlichkeit, dass ein zufällig gewählter Pendler genau eines von beiden tut: Rad oder U-Bahn, aber nicht beides}
\erg{13}{a) Tom rechnet das einschließende Oder; bei „entweder – oder“ wird der Schnitt zweimal abgezogen: $0{,}4 + 0{,}25 - 2 \cdot 0{,}1 = 0{,}45$ \quad b) In $P(A) + P(B)$ steckt der Schnitt doppelt. Für „A oder B“ soll er einmal zählen, also einmal abziehen; für „genau eines“ soll er gar nicht zählen, also zweimal abziehen. \quad c) $100\,\% - (62\,\% + 48\,\% - 25\,\%) = 15\,\%$ \quad d) $R \cap \bar{B}$}
